\documentclass[a4paper,11pt]{article}

\usepackage[a4paper,margin=2cm]{geometry}
\usepackage[T1]{fontenc}
\usepackage[utf8]{inputenc}
\usepackage[ngerman,english]{babel}
\usepackage{lmodern}
\usepackage{microtype}
\usepackage{xcolor}
\usepackage{parskip}
\usepackage{enumitem}
\usepackage[most]{tcolorbox}
\usepackage{titlesec}
\usepackage[hidelinks]{hyperref}
\usepackage{array}
\usepackage{longtable}
\usepackage[table]{xcolor}

\definecolor{HeaderBlueDark}{RGB}{70,110,170}
\definecolor{RowBlueLight}{RGB}{235,243,255}
\definecolor{RowBlueDark}{RGB}{220,233,250}
\definecolor{SoftGray}{RGB}{90,90,90}

\setlength{\parindent}{0pt}
\setlist[itemize]{leftmargin=1.4em,itemsep=0.35em,topsep=0.35em}
\linespread{1.06}
\renewcommand{\arraystretch}{1.35}
\setlength{\tabcolsep}{6pt}

\titleformat{\section}[block]
{\Large\bfseries\color{HeaderBlueDark}}
{}{0pt}{}
\titlespacing{\section}{0pt}{1.3em}{0.8em}

\titleformat{\subsection}[block]
{\large\bfseries\color{HeaderBlueDark}}
{}{0pt}{}
\titlespacing{\subsection}{0pt}{1.0em}{0.6em}

\newtcolorbox{exbox}{enhanced,breakable,colback=RowBlueLight,colframe=RowBlueDark,boxrule=0.4pt,arc=1.5mm,left=1.2mm,right=1.2mm,top=1mm,bottom=1mm,title={Beispiele \textnormal{(Examples)}},fonttitle=\bfseries,colbacktitle=RowBlueDark,coltitle=black}

\NewDocumentEnvironment{grammarentry}{mm+b}{%
    \begin{tcolorbox}[enhanced,breakable,colback=white,colframe=HeaderBlueDark,boxrule=0.6pt,arc=1.5mm,left=1.5mm,right=1.5mm,top=1mm,bottom=1mm,colbacktitle=HeaderBlueDark,coltitle=white,fonttitle=\bfseries,title={#1\hfill\normalfont\footnotesize #2}]
    #3
    \end{tcolorbox}
}{}

\newcommand{\GermanText}[1]{\textbf{Deutsch: }#1\par}
\newcommand{\EnglishText}[1]{{\small\color{SoftGray}\textbf{English: }#1\par}}
\newcommand{\ExampleItem}[2]{\item \textbf{#1}\\{\small\color{SoftGray}#2}}

\title{\glqq nach + Dativ\grqq\\in der Bedeutung \glqq gemäß / entsprechend\grqq}
\date{}

\begin{document}
\maketitle
\tableofcontents
\newpage

\section{Genau die grammatische Rolle in deinem Satz}

\begin{grammarentry}{nach den obigen Argumenten}{Präposition + Dativ}
\GermanText{In \glqq nach den obigen Argumenten\grqq ist \textbf{nach} eine \textbf{Präposition}. In genau dieser Verwendung bedeutet sie ungefähr \textbf{gemäß}, \textbf{entsprechend} oder \textbf{nach Maßgabe von}.}
\EnglishText{In \textit{nach den obigen Argumenten}, \textbf{nach} is a \textbf{preposition}. In this specific use it means approximately \textbf{according to}, \textbf{in accordance with}, or \textbf{judging by}.}

\GermanText{Die Präposition \textbf{nach} regiert hier immer den \textbf{Dativ}. Deshalb heißt es: \textbf{nach den obigen Argumenten}.}
\EnglishText{The preposition \textbf{nach} governs the \textbf{dative} here. Therefore the form is: \textbf{nach den obigen Argumenten}.}
\end{grammarentry}

\section{Wortart, Kasus und Satzfunktion}

\begin{grammarentry}{Drei Ebenen}{nicht miteinander verwechseln}
\GermanText{\textbf{Wortart von \glqq nach\grqq:} Präposition.}
\EnglishText{\textbf{Word class of \textit{nach}:} preposition.}

\GermanText{\textbf{Kasusrektion:} \textbf{nach + Dativ}. Das Nomen oder Pronomen nach \textbf{nach} steht im Dativ.}
\EnglishText{\textbf{Case government:} \textbf{nach + dative}. The noun or pronoun after \textbf{nach} is in the dative.}

\GermanText{\textbf{Satzfunktion der ganzen Gruppe:} Die Präpositionalgruppe \glqq nach den obigen Argumenten\grqq funktioniert adverbial und gibt die \textbf{Grundlage / den Bezugsrahmen einer Beurteilung oder Aussage} an.}
\EnglishText{\textbf{Syntactic function of the whole phrase:} The prepositional phrase \textit{nach den obigen Argumenten} functions adverbially and indicates the \textbf{basis / frame of reference for a judgment or statement}.}
\end{grammarentry}

\section{Warum steht \glqq den\grqq?}

\rowcolors{2}{RowBlueLight}{RowBlueDark}
\begin{longtable}{>{\bfseries}p{3.2cm}|p{4.4cm}|p{6.7cm}}
\rowcolor{HeaderBlueDark}
\color{white}\textbf{Element} &
\color{white}\textbf{Form} &
\color{white}\textbf{Erklärung} \\
Präposition & nach & verlangt in dieser Verwendung den Dativ \\
Artikel & den & bestimmter Artikel, Dativ Plural \\
Adjektiv & obigen & schwache Adjektivendung im Dativ Plural nach bestimmtem Artikel \\
Nomen & Argumenten & Dativ Plural von \glqq die Argumente\grqq; zusätzliches \textbf{-n} im Dativ Plural \\
\end{longtable}

\begin{grammarentry}{Schema}{Formbildung}
\GermanText{\textbf{nach + den + obigen + Argumenten} = Präposition + bestimmter Artikel im Dativ Plural + schwach dekliniertes Adjektiv + Nomen im Dativ Plural.}
\EnglishText{\textbf{nach + den + obigen + Argumenten} = preposition + definite article in dative plural + weakly declined adjective + noun in dative plural.}
\end{grammarentry}

\section{Bedeutung dieser Konstruktion}

\begin{grammarentry}{Bezug oder Maßstab}{according to}
\GermanText{Diese Verwendung von \textbf{nach} sagt, dass eine Aussage, Einschätzung oder Beurteilung sich an einer Quelle, Regel, Information, Meinung oder Grundlage orientiert.}
\EnglishText{This use of \textbf{nach} says that a statement, assessment, or judgment is based on or oriented toward a source, rule, piece of information, opinion, or reference basis.}

\GermanText{Man kann oft mit \textbf{gemäß}, \textbf{entsprechend} oder im Englischen mit \textbf{according to} paraphrasieren.}
\EnglishText{It can often be paraphrased with \textbf{gemäß}, \textbf{entsprechend}, or in English with \textbf{according to}.}
\end{grammarentry}

\begin{exbox}
\begin{itemize}
\ExampleItem{Nach diesen Informationen ist die Straße gesperrt.}{According to this information, the road is closed.}
\ExampleItem{Nach seiner Aussage war niemand im Raum.}{According to his statement, nobody was in the room.}
\ExampleItem{Nach den vorliegenden Daten ist die Methode erfolgreich.}{According to the available data, the method is successful.}
\ExampleItem{Nach den obigen Argumenten ist diese Schlussfolgerung plausibel.}{According to the above arguments, this conclusion is plausible.}
\end{itemize}
\end{exbox}

\section{Stellung im Satz}

\begin{grammarentry}{Adverbiale Präpositionalgruppe}{Vorfeld oder Mittelfeld}
\GermanText{Die ganze Präpositionalgruppe kann im \textbf{Vorfeld} stehen. Dann folgt im Hauptsatz sofort das finite Verb.}
\EnglishText{The whole prepositional phrase can occupy the \textbf{prefield}. In a main clause, the finite verb then follows immediately.}

\GermanText{Schema: \textbf{Nach diesen Informationen + Verb + Subjekt + ...}}
\EnglishText{Pattern: \textbf{According to this information + verb + subject + ...}}

\GermanText{Sie kann auch im \textbf{Mittelfeld} stehen.}
\EnglishText{It can also appear in the \textbf{middle field}.}
\end{grammarentry}

\begin{exbox}
\begin{itemize}
\ExampleItem{Nach diesen Informationen ist die Straße gesperrt.}{According to this information, the road is closed.}
\ExampleItem{Die Straße ist nach diesen Informationen gesperrt.}{According to this information, the road is closed.}
\end{itemize}
\end{exbox}

\section{Was \glqq nach\grqq hier nicht ausdrückt}

\begin{grammarentry}{Nur diese eine Rolle}{Abgrenzung}
\GermanText{In dieser Datei geht es \textbf{nicht} um \textbf{nach} als Zeitangabe wie in \glqq nach dem Essen\grqq und \textbf{nicht} um eine Richtung wie in \glqq nach Wien fahren\grqq.}
\EnglishText{This file is \textbf{not} about \textbf{nach} as a temporal expression such as \textit{nach dem Essen}, and \textbf{not} about direction such as \textit{nach Wien fahren}.}

\GermanText{Hier wird ausschließlich die Verwendung behandelt, in der \textbf{nach + Dativ} einen \textbf{Bezugsrahmen, eine Quelle oder einen Maßstab für eine Aussage} nennt.}
\EnglishText{Only the use in which \textbf{nach + dative} names a \textbf{frame of reference, source, or standard for a statement} is treated here.}
\end{grammarentry}

\section{Unterschied zu \glqq aus den Argumenten\grqq}

\begin{grammarentry}{nach = Bezugsrahmen}{nicht Herkunft einer Folgerung}
\GermanText{\textbf{Nach den Argumenten} präsentiert die Argumente als \textbf{Bezugsrahmen für die folgende Aussage}.}
\EnglishText{\textbf{Nach den Argumenten} presents the arguments as the \textbf{frame of reference for the following statement}.}

\GermanText{\textbf{Aus den Argumenten} präsentiert die Argumente dagegen als \textbf{Ausgangspunkt oder Quelle, aus der eine Folgerung entsteht}. Deshalb steht \textbf{aus} besonders typisch mit Verben wie \textbf{sich ergeben}, \textbf{folgen} oder \textbf{schließen}.}
\EnglishText{\textbf{Aus den Argumenten}, by contrast, presents the arguments as the \textbf{starting point or source from which a conclusion is derived}. Therefore \textbf{aus} is especially typical with verbs such as \textbf{sich ergeben}, \textbf{folgen}, or \textbf{schließen}.}
\end{grammarentry}

\section{Typische Fehler}

\begin{grammarentry}{Fehler vermeiden}{B1-relevant}
\GermanText{Falsch: \textbf{nach die Argumente}. Richtig: \textbf{nach den Argumenten}. Grund: \textbf{nach} verlangt hier den Dativ.}
\EnglishText{Wrong: \textbf{nach die Argumente}. Correct: \textbf{nach den Argumenten}. Reason: \textbf{nach} requires the dative here.}

\GermanText{Falsch: \textbf{nach den obige Argumenten}. Richtig: \textbf{nach den obigen Argumenten}. Grund: Nach \textbf{den} bekommt das Adjektiv im Dativ Plural die schwache Endung \textbf{-en}.}
\EnglishText{Wrong: \textbf{nach den obige Argumenten}. Correct: \textbf{nach den obigen Argumenten}. Reason: after \textbf{den}, the adjective takes the weak ending \textbf{-en} in dative plural.}
\end{grammarentry}

\section{Merksatz}

\begin{grammarentry}{Kurzform}{für deinen Satz}
\GermanText{\textbf{nach + Dativ} = hier: \textbf{gemäß / entsprechend / according to}. Die ganze Präpositionalgruppe gibt den \textbf{Bezugsrahmen der Aussage} an.}
\EnglishText{\textbf{nach + dative} = here: \textbf{according to / in accordance with}. The whole prepositional phrase gives the \textbf{frame of reference for the statement}.}
\end{grammarentry}

\end{document}
